\section{Conclusion}
\label{sec:conclusion}

We presented \textbf{SETA}, the strongest frozen-feature recipe for
DAiSEE engagement recognition in the literature: bagged class-balanced
logistic-regression probes on SigLIP-L with val-tuned ordinal
threshold decoding, achieving $\kappa_q = 0.247~[0.201, 0.294]$ ---
a 24\% relative improvement over the standard linear-probe argmax
baseline, with statistically meaningful lower-CI separation
(0.201 $>$ 0.199). The recipe is single-encoder, parameter-light,
single-CPU reproducible, and provides a calibrated baseline for any
subsequent work on this benchmark.

We further introduced two novel diagnostics --- the
\emph{identity-engagement entanglement} (IDEP) curve and the
\emph{within-subject ranking} analysis --- that together explain
\emph{why} SETA is the achievable frozen-feature ceiling. Identity
is encoded across dozens of redundant directions while engagement
variance is concentrated in only two-to-five linear directions that
overlap with identity; any linear identity-erasure also erases
engagement (LEACE collapses $\kappa_q$ to $0.04$). Within-subject
Spearman $\rho \approx 0.05$: even at SETA's global $\kappa_q$, the
model captures subject-level priors and not per-clip discrimination.
DREAM, our self-supervised neutral-anchor personalization scheme,
constructively falsifies the ``remove identity, recover
engagement'' hypothesis at the readout level.

Our systematic decomposition across eight encoder/data-adaptation
families and sixty-plus method variants demonstrates that the
frozen-feature regime cleanly plateaus at $\kappa_q \in [0.20, 0.25]$.
The cross-task control --- the same frozen encoders reach
$\kappa = 0.55$ on FER2013 emotion classification --- bounds the
claim: the DAiSEE ceiling is engagement-task-specific, not
encoder-generic. The path forward, based on our diagnostic, is
end-to-end engagement-supervised retraining (which the existing
supervised SOTA ViBED-Net achieves at $\kappa \approx 0.55$ via 3-D
CNN end-to-end training); we release a GPU-ready script
(ENGAGENET-X, supplementary) as a reproducible starting point for
the next paper in this line.

SETA fills the missing positive baseline for frozen-feature DAiSEE
engagement. The IDEP curve and within-subject ranking analysis are
the diagnostic tools that explain the gap to supervised SOTA. We
hope this combined recipe-and-decomposition becomes the canonical
reference for the next generation of work on classroom engagement
recognition.
